\section{Privacy, Encryption y señales mínimas necesarias}

Los sistemas de seguridad procesan contenido muy sensible e información de cuentas. La privacidad no es una «restricción externa» opuesta a la seguridad, sino un objetivo de arquitectura. El \term{end-to-end encryption} (cifrado de extremo a extremo, E2EE) significa que solo los extremos de la comunicación tienen las claves del contenido de los mensajes; ni el transporte ni el servidor pueden leer el texto en claro. Protege a los usuarios de que la plataforma, los atacantes y los intermediarios espíen el contenido, y también limita el content scanning del lado del servidor. Un diseño responsable debe precisar qué señales siguen disponibles, quién puede acceder a ellas y cómo evitar que la vigilancia se amplíe.

\subsection{Cuatro capas: content, metadata, endpoint y governance}

Esta sección no busca una solución universal para «eludir el cifrado», sino que analiza el problema por capas. Los \term{metadata} (metadatos) son las propiedades de la comunicación o del sistema distintas del contenido, por ejemplo la hora, la cuenta, el dispositivo, las relaciones de conexión y los eventos del servicio; pueden ser muy sensibles para la privacidad, y no deben recopilarse sin límite solo porque no forman parte del cuerpo del mensaje. Los endpoint controls se sitúan en el dispositivo del usuario o en el límite de la aplicación; las network signals describen comportamientos coordinados, y la governance restringe el acceso, el legal process, la appeal y la oversight.

\lecturefigure{11-privacy-layers.png}{La privacidad y la seguridad infantil requieren controles por capas de content, metadata, endpoint y governance.}{Subtítulos locales 28:40--31:10; principios de NIST y de gobernanza de plataformas; redibujo conceptual.}

\begin{knowledgebox}{Lectura de la figura: cada señal debe declarar su purpose limitation}
La content signal puede ser la más directa, pero también la más invasiva; los metadata pueden revelar relaciones anómalas, pero es fácil que afecten por error conductas sociales ordinarias; la endpoint intervention puede alertar antes al usuario, pero hay que evitar que se amplíe la vigilancia a nivel de dispositivo; la governance determina quién puede consultar, cuándo se elimina la información y cómo se apela. Para cada capa deben especificarse el purpose, la base legal, los campos mínimos, la retention, el access role, la audit y los secondary use prohibidos.
\end{knowledgebox}

\begin{warningbox}{«La privacidad y la seguridad son importantes» no es una respuesta de diseño}
Una respuesta de ingeniería real necesita un threat model y límites: si se altera la confidencialidad de extremo a extremo, si se introduce una capacidad de escaneo que pueda reutilizarse con otros fines, quién corrige los falsos positivos, si los menores de edad y los grupos vulnerables asumen riesgos adicionales, y si el gobierno o el personal interno pueden ampliar las consultas. Cualquier propuesta debe someterse a una evaluación de seguridad independiente, a una evaluación de impacto en los derechos y a una gobernanza transparente, en lugar de depender solo de las promesas del proveedor.
\end{warningbox}

\teachervoice{Cordua reconoce que el encryption impone limitaciones reales, y a la vez subraya que la plataforma debe seguir buscando vías de seguridad legítimas y que protejan la privacidad. Lo más valioso de la clase no es una postura política concreta, sino la idea de «no fingir que el trade-off no existe, y tampoco dejar por ello de diseñar el sistema».}

\subsection{Resumen de la sección}

La privacidad y la seguridad requieren mínimo privilegio, purpose limitation y una gobernanza auditable. El E2EE protege la confidencialidad del contenido y cambia lo que el servidor puede ver; los metadata, las endpoint signals o las network signals solo pueden usarse como señales de riesgo por capas, y no se convierten automáticamente en justificación de una vigilancia amplia.
